\paragraph{伴随矩阵、Witt 分解、模 \texorpdfstring{$m$}{m} 可解性与交叠块的整数逆核}
沿用 \S\ref{subsubsec:xi-time-protocol-conclusions-part30a-boolean-kernel-rigidity} 的记号：
固定 \(q\ge 2\)，令 \(V_q:=2^{[q]}\) 并以集合交并书写指标；
记 \(A_q:=2T^{(q)}\in\Mat_{2^q}(\ZZ)\) 由
\[
(A_q)_{U,V}=
\begin{cases}
2,&U\cap V=\varnothing,\\
1,&U\cap V\neq \varnothing,
\end{cases}
\qquad (U,V\subseteq[q])
\]
给出。

\begin{corollary}[伴随矩阵的覆盖支撑与带符号余因子值域]\label{cor:xi-disjointness-boolean-adjugate-cover-support}
对任意 \(U,V\in V_q\) 有
\[
\operatorname{adj}(A_q)_{U,V}
=
2(-1)^{\card{U\cap V}}\mathbf 1_{\{U\cup V=[q]\}}
-\mathbf 1_{\{U=\varnothing=V\}}.
\]
因此任意 \((2^q-1)\times(2^q-1)\) 的带符号余因子仅取值于 \(\{0,-1,\pm2\}\)；
并且除 \((U,V)=(\varnothing,\varnothing)\) 外，\(\operatorname{adj}(A_q)_{U,V}\neq 0\) 当且仅当 \(U\cup V=[q]\)。
\end{corollary}

\begin{proof}
由结论 \ref{con:xi-terminal-disjointness-boolean-det-rigidity} 知 \(\det(A_q)=2\)，并由结论 \ref{con:xi-terminal-disjointness-boolean-inverse-closed-form} 知
\(\operatorname{adj}(A_q)=\bigl(T^{(q)}\bigr)^{-1}\)。
将其中的逐元闭式从按位 \((\wedge,\vee)\) 记号翻译为集合 \((\cap,\cup)\) 记号（并注意 \(\mathbf 0\leftrightarrow\varnothing\)、\(\mathbf 1\leftrightarrow[q]\)）即得所示公式。
余因子值域与非零支撑随之成立。证毕。
\end{proof}

\begin{proposition}[有理二次型的 Witt 分解与各向异核的稳定性]\label{prop:xi-disjointness-boolean-witt-decomposition-q}
视 \(A_q\) 为 \(\QQ\)-二次型的 Gram 矩阵，并在 \(\QQ^{V_q}\) 上定义
\[
Q_q(x):=x^{\mathsf T}A_qx,\qquad x\in\QQ^{V_q}.
\]
则对一切 \(q\ge 2\) 成立正交分解
\[
(\QQ^{V_q},Q_q)\ \cong\ H^{\oplus (2^{q-1}-1)}\ \perp\ \langle 1,-2\rangle,
\]
其中 \(H\) 为双曲平面。
因而 \(Q_q\) 的 Witt 指数为 \(2^{q-1}-1\)，并且其各向异核在 \(\QQ\) 上同构于 \(\langle 1,-2\rangle\)，与 \(q\) 无关。
\end{proposition}

\begin{proof}
由命题 \ref{prop:xi-disjointness-boolean-zeta-ldlt} 的整同余对角化
\(\mu_q^{\mathsf T}A_q\mu_q=\Delta_q\)，可知 \(Q_q\) 在 \(\QQ\) 上同构于对角型
\[
2x_{\varnothing}^2+\sum_{\substack{S\neq\varnothing\\ \card{S}\ \mathrm{even}}}x_S^2
-\sum_{\substack{S\\ \card{S}\ \mathrm{odd}}}x_S^2.
\]
非空偶子集数为 \(2^{q-1}-1\)，奇子集数为 \(2^{q-1}\)。
将每个非空偶子集与一个奇子集配对，可逐对剥离出双曲平面 \(\langle 1\rangle\perp\langle-1\rangle\cong H\)，共得到 \(2^{q-1}-1\) 个 \(H\)。
剩余的二维部分同构于 \(\langle 2,-1\rangle\)，经尺度变换等价于 \(\langle 1,-2\rangle\)。
最后，若 \(\langle 1,-2\rangle\) 在 \(\QQ\) 上各向同性，则存在 \((a,b)\neq (0,0)\) 使 \(a^2-2b^2=0\)，从而 \(2=(a/b)^2\) 为有理平方，矛盾。
故该二维部分为各向异核，Witt 分解成立。证毕。
\end{proof}

\begin{corollary}[\texorpdfstring{$\sqrt2$}{sqrt2} 判据与 \texorpdfstring{$p$}{p}-进分型]\label{cor:xi-disjointness-boolean-witt-sqrt2-criterion}
令 \(F\) 为特征不为 \(2\) 的域，并将 \(Q_q\) 标量扩张为 \(Q_{q,F}\)。
则对一切 \(q\ge 2\) 有
\[
(F^{V_q},Q_{q,F})\ \cong\ H^{\oplus (2^{q-1}-1)}\ \perp\ \langle 1,-2\rangle_F.
\]
从而该二次空间在 \(F\) 上全分裂为双曲型（即同构于 \(H^{\oplus 2^{q-1}}\)）当且仅当 \(2\in(F^\times)^2\)。
特别地，当 \(F=\QQ_p\)（\(p\) 为奇素数）时，
\[
(\QQ_p^{V_q},Q_{q,\QQ_p})\ \text{全分裂}
\quad\Longleftrightarrow\quad
\left(\frac{2}{p}\right)=1
\quad\Longleftrightarrow\quad
p\equiv \pm 1\pmod 8.
\]
\end{corollary}

\begin{proof}
命题 \ref{prop:xi-disjointness-boolean-witt-decomposition-q} 的正交直和分解在任意标量扩张下保持，得第一式。
二维型 \(\langle 1,-2\rangle_F\) 为双曲平面当且仅当存在非零解使 \(x^2-2y^2=0\)，等价于 \(2\) 为 \(F\) 中平方。
对 \(F=\QQ_p\)（\(p\) 奇），\(\sqrt2\in\QQ_p\) 当且仅当 \(2\) 为二次剩余；由二次互反律的标准结论
\(\left(\frac{2}{p}\right)=(-1)^{(p^2-1)/8}\)，得到模 \(8\) 的等价刻画。证毕。
\end{proof}

\begin{proposition}[模 \texorpdfstring{$m$}{m} 线性方程的存在性与解数：\texorpdfstring{$\gcd(2,m)$}{gcd(2,m)} 规律]\label{prop:xi-disjointness-boolean-congruence-solvability-gcd2m}
令 \(m\ge 1\)。
在 \((\ZZ/m\ZZ)^{V_q}\) 上考虑同余方程
\[
A_qx\equiv b\pmod m.
\]
则对一切 \(q\ge 2\)，该方程可解当且仅当
\[
b_{\varnothing}\equiv 0\pmod{\gcd(2,m)}.
\]
当可解时，其解集为一个仿射陪集，解的总数恰为 \(\gcd(2,m)\)。
更精确地：
\begin{enumerate}
  \item 若 \(m\) 为奇数，则解唯一；
  \item 若 \(m\) 为偶数，则恰有两解，且其差为 \((m/2)\delta_{\varnothing}\)。
\end{enumerate}
\end{proposition}

\begin{proof}
由结论 \ref{con:xi-terminal-disjointness-boolean-smith-mod2-rank-kernel} 的 Smith 形，
\(
\mathrm{SNF}(A_q)=\mathrm{diag}(1,\dots,1,2)
\)；
从而在任意模 \(m\) 下，方程可解的唯一障碍来自唯一的因子 \(2\)，即首坐标必须落入 \(\gcd(2,m)\) 倍子群。
另一方面，由命题 \ref{prop:xi-disjointness-boolean-zeta-ldlt} 的整同余分解 \(A_q=Z_q^{\mathsf T}\Delta_q Z_q\)，取 \(y:=Z_qx\)、\(c:=Z_q^{-\mathsf T}b\)，可将同余方程化为
\(\Delta_q y\equiv c\pmod m\)；
其中除 \(\varnothing\) 分量外的对角元均为 \(\pm 1\)，故唯一非平凡约束为
\(
2y_{\varnothing}\equiv c_{\varnothing}\pmod m
\)。
又 Möbius 矩阵满足 \((Z_q^{-\mathsf T}b)_{\varnothing}=b_{\varnothing}\)，因此可解条件等价于 \(b_{\varnothing}\equiv 0\pmod{\gcd(2,m)}\)。
当该条件满足时，\(\varnothing\) 分量的解数正好为 \(\gcd(2,m)\)，其余分量唯一确定，故总解数亦为 \(\gcd(2,m)\)。
偶模情形下核由 \((m/2)\delta_{\varnothing}\) 生成，给出“两解差”为 \((m/2)\delta_{\varnothing}\) 的精确描述。证毕。
\end{proof}

\begin{proposition}[非空交叠矩阵的整数逆核：支撑在覆盖对且仅取 \texorpdfstring{$\{-1,0,1\}$}{-1,0,1}]\label{prop:xi-disjointness-boolean-overlap-matrix-inverse}
令 \(V_q^+:=V_q\setminus\{\varnothing\}\)，并定义交叠矩阵 \(B_q\in \ZZ^{V_q^+\times V_q^+}\)：
\[
(B_q)_{U,V}:=\mathbf 1_{\{U\cap V\neq\varnothing\}}\qquad(U,V\in V_q^+).
\]
则 \(B_q\) 在 \(\ZZ\) 上可逆，且其逆矩阵满足
\[
\bigl(B_q\bigr)^{-1}_{U,V}
=(-1)^{\card{U\cap V}+1}\mathbf 1_{\{U\cup V=[q]\}}
\qquad (U,V\in V_q^+).
\]
因此 \(\bigl(B_q\bigr)^{-1}\) 的每个元素都属于 \(\{-1,0,1\}\)，并且其非零位置精确由覆盖条件 \(U\cup V=[q]\) 刻画。
\end{proposition}

\begin{proof}
取基变换 \(P\in\mathrm{GL}_{2^q}(\ZZ)\) 由
\[
f_{\varnothing}=\delta_{\varnothing},\qquad
f_U=\delta_U-\delta_{\varnothing}\ \ (U\in V_q^+)
\]
定义。
直接计算得到块分解
\[
P^{\mathsf T}A_qP=\begin{pmatrix}2&0\\[2pt]0&-B_q\end{pmatrix}.
\]
两边取逆并比较右下块，得到
\[
\bigl(A_q^{-1}\bigr)\big|_{V_q^+\times V_q^+}=-(B_q)^{-1}.
\]
另一方面，由结论 \ref{con:xi-terminal-disjointness-boolean-inverse-closed-form} 与 \(\operatorname{adj}(A_q)=2A_q^{-1}\) 可知：
对 \(U,V\in V_q^+\)，有
\[
2(A_q^{-1})_{U,V}
=\operatorname{adj}(A_q)_{U,V}
=2(-1)^{\card{U\cap V}}\mathbf 1_{\{U\cup V=[q]\}},
\]
从而 \((A_q^{-1})_{U,V}=(-1)^{\card{U\cap V}}\mathbf 1_{\{U\cup V=[q]\}}\)。
代回即可得到 \((B_q)^{-1}\) 的逐元闭式。证毕。
\end{proof}

\endinput

